\newpage
\renewcommand{\theequation}{\Alph{section}\arabic{equation}}
\setcounter{equation}{0}
\appendix
% \addappheadtotoc

\chapter{Codis}\label{ap.A}
Per comprovar els autovalors s'ha fet servir una rutina de Mathematica que els calcula amb precisió i determina el seu comportament.

Totes les solucions a les equacions, càlculs i gràfiques presentats en aquest treball s'han realitzat en notebooks
de Python (format .ipynb), amb l'ajuda de la llibreria Dedalus.
Per millorar l'eficiència i la qualitat del codi, 
s'ha fet servir l'assistent d'intel·ligència artificial de Copilot durant el desenvolupament del projecte.
Aquest està implementat en VS Code, l'eina amb la que s'han realitzat els notebooks i amb la que s'ha redactat la present memòria.

Tant el quadern de Mathematica com els codis de Python es troben en el repositori de GitHub \parencite{github}.


\chapter{Valors realistes} \label{ap.B}
En aquest apèndix es mostren els resultats obtinguts emprant Dedalus per valors realistes de les velocitats del so i d'Alfvén de les Eqs.~(\ref{eq.3.34})~i~(\ref{eq.3.35}).
Llavors aquests substitueixen els de la Taula~\ref{tab:valors} pels de la Taula~\ref{tab:valors_apendix}.

\begin{table}[H]
	\centering
	\begin{tabular}{|c|c|c|c|c|}
		\hline
		$c_{s c}^2$ & $c_{s p}^2$ & $v_{A c}^2$ & $v_{A p}^2$ & $k_z$ \\
		\hline
		$12$ & $166$ & $60$ & $1400$ & $0.01$ \\
		\hline
	\end{tabular}
	\caption{Valors numèrics de velocitats adimensionals realistes i constant d'acoblament del sistema.}
	\label{tab:valors_apendix}
\end{table}
\vspace{-10pt}

\begin{figure}[H]
    \centering
    \includegraphics[width=0.95\textwidth]{Figures_F&S_modes/eigenvalue_error_apendix.png}
	\caption{Error relatiu dels autovalors en el cas dels modes ràpids i lents.
	Els autovalors calculats amb Dedalus, $\omega_D$, estan calculats pel cas $k_z = 0.01$,
    llavors això afecta negativament al càlcul de l'error.}
	\label{fig:error_relatiu_apendix}
\end{figure}



\chapter{Imatges} \label{ap.C}
En aquest apèndix s'introdueixen les imatges que no s'han presentat en el cos de la memòria per qüestions d'espai.

\begin{figure}[H]
    \centering
    \includegraphics[width=0.95\textwidth]{Representacions/Solar_prominence_wiki.png}
	\caption{Protuberància real. Imatge obtinguda de la Wikipèdia.
    \url{https://en.wikipedia.org/wiki/File:Solar_prominence_2011-04-14T202956.120.png}}
	\label{fig:real_prominence}
\end{figure}

\begin{figure}[H]
    \centering
    \includegraphics[width=0.95\textwidth]{Representacions/faq1.jpg}
	\caption{Tamany aproximat de la Terra comparat amb el d'una protuberància real.
    Imatge treobtinguda de la NASA.
    \url{https://www.nasa.gov/image-article/what-solar-prominence/}}
	\label{fig:real_prominence_earth}
\end{figure}


